\documentclass[11pt,a4paper]{article}

% -----------------------------------------------------------------------------
% Long Assessment 1 report skeleton
% Adapt, combine, or reorder sections to support your own scientific argument.
% Delete all TODO text and unused examples before submission.
% -----------------------------------------------------------------------------

\usepackage[margin=25mm]{geometry}
\usepackage[T1]{fontenc}
\usepackage{amsmath,amssymb,bm}
\usepackage{booktabs,longtable,tabularx,array,multirow}
\usepackage{graphicx}
\usepackage{siunitx}
\usepackage[dvipsnames]{xcolor}
\usepackage{enumitem}
\usepackage{fancyhdr}
\usepackage{microtype}
\usepackage{caption}
\usepackage{subcaption}
\usepackage{hyperref}
\usepackage[nameinlink,noabbrev]{cleveref}

\hypersetup{
  colorlinks=true,
  linkcolor=MidnightBlue,
  citecolor=ForestGreen,
  urlcolor=BrickRed,
  pdftitle={SEMPHY01 Long Assessment 1},
  pdfauthor={Student Name}
}

\sisetup{
  scientific-notation=true,
  separate-uncertainty=true,
  per-mode=symbol
}

\graphicspath{{../figures/}{figures/}}
\setlist{nosep}
\setlength{\parindent}{0pt}
\setlength{\parskip}{0.65em}
\renewcommand{\arraystretch}{1.15}

\pagestyle{fancy}
\setlength{\headheight}{14pt}
\fancyhf{}
\lhead{SEMPHY01 --- Long Assessment 1}
\rhead{\studentname}
\cfoot{\thepage}

% ---------- Student-editable metadata ----------------------------------------
\newcommand{\studentname}{TODO: Student Name}
\newcommand{\studentid}{TODO: Student ID}
\newcommand{\materialname}{TODO: Material and phase/polytype}
\newcommand{\submissiondate}{TODO: Submission date}
\newcommand{\repositoryurl}{https://example.invalid/replace-me}
\newcommand{\physkitversion}{TODO: PhysKit version/commit}

% ---------- Frequently used notation ----------------------------------------
\newcommand{\kB}{k_{\mathrm B}}
\newcommand{\Ec}{E_{\mathrm C}}
\newcommand{\Ev}{E_{\mathrm V}}
\newcommand{\Eg}{E_{\mathrm g}}
\newcommand{\mui}{\mu_{\mathrm i}}
\newcommand{\Nc}{N_{\mathrm C}}
\newcommand{\Nv}{N_{\mathrm V}}
\newcommand{\nIntrinsic}{n_{\mathrm i}}
\newcommand{\todo}[1]{\textcolor{BrickRed}{\textbf{TODO:} #1}}

% A compiling placeholder for figures that are not ready yet.
\newcommand{\placeholderfigure}[3][0.82\linewidth]{%
  \begin{figure}[htbp]
    \centering
    \fbox{\parbox[c][45mm][c]{#1}{\centering #2}}
    \caption{#2}
    \label{#3}
  \end{figure}
}

\begin{document}

\begin{titlepage}
  \centering
  {\Large De La Salle University\par}
  \vspace{8mm}
  {\large SEMPHY01 --- Semiconductor Physics\par}
  \vspace{18mm}
  {\Huge\bfseries Long Assessment 1\par}
  \vspace{4mm}
  {\LARGE From Bloch Bands to Equilibrium Semiconductor Carrier Statistics\par}
  \vspace{18mm}

  \begin{tabular}{@{}rl@{}}
    Student: & \studentname \\
    Student ID: & \studentid \\
    Selected material: & \materialname \\
    Submission date: & \submissiondate \\
    Computational record: & \url{\repositoryurl} \\
    PhysKit version: & \physkitversion
  \end{tabular}

  \vfill
  \begin{minipage}{0.88\textwidth}
    \small
    This report concerns equilibrium carrier statistics only. Transport,
    mobility, conductivity, device behavior, and spatial band bending are
    outside the model unless explicitly identified as excluded context.
  \end{minipage}
\end{titlepage}

\pagenumbering{roman}
\tableofcontents
\clearpage
\pagenumbering{arabic}

% =============================================================================
\section{Project declaration and scope}
\label{sec:declaration}
% =============================================================================

\begin{table}[htbp]
  \centering
  \caption{Declared scope of the material model.}
  \label{tab:declaration}
  \begin{tabularx}{\textwidth}{@{}>{\bfseries}p{0.25\textwidth}X@{}}
    \toprule
    Item & Declaration \\
    \midrule
    Material identity & \todo{Chemical formula and exact phase/polytype.} \\
    Sample model & \todo{Ideal bulk model and temperature range.} \\
    Gap model & \todo{Direct/indirect character and $\Eg(T)$ representation.} \\
    Conduction edge & \todo{Location, equivalent valleys, and mass model.} \\
    Valence edge & \todo{Relevant bands, degeneracies, and mass model.} \\
    Statistics & \todo{Full Fermi--Dirac, Boltzmann comparison, or both.} \\
    n-type model & \todo{Donor-like species/defect and charge states.} \\
    p-type model & \todo{Acceptor-like species/defect and charge states.} \\
    Exclusions & \todo{Omitted physics and why it is outside scope.} \\
    \bottomrule
  \end{tabularx}
\end{table}

\subsection{Research question and modeling strategy}

\todo{State the material-specific question investigated and summarize the path
from the Bloch calculation to intrinsic and doped equilibrium statistics.}

\subsection{Use of PhysKit and student-authored work}

The complete PhysKit library may be used \cite{physkit}. Distinguish maintained
package capabilities from your own configuration, analysis, extension, and
interpretation.

\begin{table}[htbp]
  \centering
  \caption{Computational capability and contribution map.}
  \label{tab:contributions}
  \begin{tabularx}{\textwidth}{@{}p{0.25\textwidth}p{0.23\textwidth}X@{}}
    \toprule
    Capability & Source & Configuration, extension, or validation performed \\
    \midrule
    Bloch Hamiltonian & PhysKit & \todo{Version, model, parameters, and checks.} \\
    DOS/state counting & \todo{Source} & \todo{Normalization and validation.} \\
    Carrier statistics & \todo{Source} & \todo{Integration/root-solving choices.} \\
    Dopant model & \todo{Source} & \todo{Charge-state convention and extension.} \\
    Visualization & \todo{Source} & \todo{How plots consume calculated results.} \\
    \bottomrule
  \end{tabularx}
\end{table}

\subsection{Model map}

\placeholderfigure{Replace with one compact diagram connecting
$H(\mathbf{k})$, $E_n(\mathbf{k})$, effective masses, DOS, occupation,
neutrality, and calculated carriers.}{fig:model-map}

% =============================================================================
\section{Part I: Bloch2D effective-mass reduction}
\label{sec:bloch2d}
% =============================================================================

\subsection{Parent model and numerical representation}

\todo{Describe the direct/reciprocal lattice, cell, boundary conditions,
potential, representation, reciprocal cutoff or discretization, and units.
Explain what PhysKit supplies and what you configured.}

A useful operator statement is
\begin{equation}
  H_{\mathbf G\mathbf G'}(\mathbf k)
  = \frac{\hbar^2\lvert\mathbf k+\mathbf G\rvert^2}{2m}
    \delta_{\mathbf G\mathbf G'} + V_{\mathbf G-\mathbf G'}.
  \label{eq:bloch-hamiltonian}
\end{equation}

\todo{Report a Hermiticity or equivalent operator-consistency check.}

\subsection{Band path and full-zone calculation}

\placeholderfigure{High-symmetry band path. Explain why it is interpretive but
not a DOS integration mesh.}{fig:band-path}

\placeholderfigure{Uniform two-dimensional Brillouin-zone edge map with located
extrema.}{fig:edge-map}

\subsection{Band edges and local quadratic model}

\begin{table}[htbp]
  \centering
  \caption{Located band edges and fit domains.}
  \label{tab:band-edges}
  \begin{tabular}{@{}llllll@{}}
    \toprule
    Edge & Band/sheet & $k_x$ & $k_y$ & Energy & Classification \\
    \midrule
    Valence maximum & \todo{ } & \todo{ } & \todo{ } & \todo{ } & \todo{ } \\
    Conduction minimum & \todo{ } & \todo{ } & \todo{ } & \todo{ } & \todo{ } \\
    \bottomrule
  \end{tabular}
\end{table}

For a conduction minimum,
\begin{equation}
  E_c(\mathbf{k}_c+\mathbf q)
  \simeq \Ec + \frac{\hbar^2}{2}
  \mathbf q^{\mathsf T}\mathbf M_c^{-1}\mathbf q,
\end{equation}
with the sign reversed before defining a positive hole-mass tensor at a
valence maximum.

\todo{Describe the multidimensional fit, energy/radius window, tensor
components, principal masses, principal directions, and any degeneracy or band
folding that complicates the fit.}

\begin{table}[htbp]
  \centering
  \caption{Effective-mass tensor results.}
  \label{tab:mass-tensors}
  \begin{tabular}{@{}lrrrrrr@{}}
    \toprule
    Edge & $(M^{-1})_{xx}$ & $(M^{-1})_{xy}$ & $(M^{-1})_{yy}$
         & $m_1/m_0$ & $m_2/m_0$ & Fit window \\
    \midrule
    Valence & \todo{ } & \todo{ } & \todo{ } & \todo{ } & \todo{ } & \todo{ } \\
    Conduction & \todo{ } & \todo{ } & \todo{ } & \todo{ } & \todo{ } & \todo{ } \\
    \bottomrule
  \end{tabular}
\end{table}

\subsection{Convergence, residuals, and range of validity}

\placeholderfigure{Fit-window mass drift, mesh/cutoff convergence, and/or
quadratic-fit residuals.}{fig:curvature-convergence}

\todo{Separate representation error, $\mathbf k$-mesh error, fit-window error,
and physical nonparabolicity. State a bounded energy range over which the
parabolic approximation is adequate.}

\subsection{Two-dimensional DOS and state counting}

For one anisotropic parabolic valley,
\begin{equation}
  m_{d,2D}=\sqrt{m_1m_2}, \qquad
  g_{2D}(E)=\frac{g_sg_vm_{d,2D}}{2\pi\hbar^2}\Theta(E-\Ec).
\end{equation}

\placeholderfigure{Numerical full-zone DOS and analytic parabolic DOS near the
selected edge.}{fig:dos-2d}

\todo{Verify integrated state counts and discuss mesh spacing, broadening,
finite bandwidth, additional valleys/sheets, and nonparabolicity. Emphasize
that these are two-dimensional sheet quantities.}

% =============================================================================
\section{Part II: Intrinsic three-dimensional model}
\label{sec:intrinsic}
% =============================================================================

\subsection{Band-edge state model and DOS components}

\todo{Define the energy reference, gap, all conduction valleys and valence
bands, principal masses, and degeneracies. Explain whether quoted DOS masses
already include valley multiplicity.}

For each three-dimensional parabolic component,
\begin{equation}
  g_c(E)=\frac{g_sg_v}{4\pi^2}
  \left(\frac{2}{\hbar^2}\right)^{3/2}
  \sqrt{\det\mathbf M_c}\sqrt{E-\Ec}\,\Theta(E-\Ec),
  \label{eq:dos-3d}
\end{equation}
with an analogous expression below $\Ev$ for holes.

\begin{table}[htbp]
  \centering
  \scriptsize
  \caption{Components of the effective densities of states at a representative temperature.}
  \label{tab:nc-nv-components}
  \begin{tabular}{@{}lrrrrrl@{}}
    \toprule
    Component & $g_s$ & $g_v$ & $m_1/m_0$ & $m_2/m_0$ & $m_3/m_0$ & Contribution \\
    \midrule
    \todo{Conduction valley} & \todo{ } & \todo{ } & \todo{ } & \todo{ } & \todo{ } & \todo{ } \\
    \todo{Valence band} & \todo{ } & \todo{ } & \todo{ } & \todo{ } & \todo{ } & \todo{ } \\
    \bottomrule
  \end{tabular}
\end{table}

\subsection{Fermi--Dirac occupation and carrier integrals}

\begin{align}
  f(E,\mu,T) &= \frac{1}{1+\exp[(E-\mu)/(\kB T)]}, \\
  n(\mu,T) &= \int_{\Ec}^{\infty}g_c(E)f(E,\mu,T)\,dE, \\
  p(\mu,T) &= \int_{-\infty}^{\Ev}g_v(E)[1-f(E,\mu,T)]\,dE.
\end{align}

\todo{Explain stable evaluation, units, integration method, cutoffs, tolerances,
and convergence evidence.}

\subsection{Intrinsic neutrality and nondegenerate comparison}

\begin{equation}
  n(\mui,T)-p(\mui,T)=0.
\end{equation}

\todo{Report the root bracket, tolerance, residual, and temperature sweep.}

\begin{align}
  \nIntrinsic &= \sqrt{\Nc\Nv}\exp\left(-\frac{\Eg}{2\kB T}\right), \\
  \mui &= \frac{\Ec+\Ev}{2}
  +\frac{\kB T}{2}\ln\left(\frac{\Nv}{\Nc}\right).
\end{align}

\placeholderfigure{Three-dimensional DOS and occupation at one temperature.}{fig:dos-3d}
\placeholderfigure{Intrinsic chemical potential and carrier concentration versus temperature.}{fig:intrinsic-temperature}

\todo{Compare full Fermi--Dirac integration with the nondegenerate limit and
identify when they agree.}

\subsection{Silicon control and validation evidence}

\begin{table}[htbp]
  \centering
  \caption{Silicon control at \SI{300}{\kelvin}.}
  \label{tab:silicon-control}
  \begin{tabular}{@{}lllll@{}}
    \toprule
    Quantity & Input/result & Reference value & Difference & Interpretation \\
    \midrule
    $\Eg$ & \todo{ } & approximately \SI{1.12}{\electronvolt} & \todo{ } & \todo{ } \\
    $\nIntrinsic$ & \todo{ } & approximately $1.5\times10^{10}\,\si{\per\centi\meter\cubed}$ & \todo{ } & \todo{ } \\
    $\mui$ & \todo{ } & \todo{source-dependent} & \todo{ } & \todo{ } \\
    \bottomrule
  \end{tabular}
\end{table}

\todo{Summarize automated or reproducible checks for monotonicity, energy-shift
invariance, isotropic/anisotropic limits, $T^{3/2}$ scaling, neutrality, and
integration convergence.}

% =============================================================================
\section{Part III: Literature-based material reconstruction}
\label{sec:material}
% =============================================================================

\subsection{Source strategy and parameter provenance}

\todo{Explain how curated sources, primary literature, reviews, and published
band information were reconciled. Do not mix phases or definitions silently.}

\scriptsize
\begin{longtable}{@{}p{0.14\textwidth}p{0.08\textwidth}p{0.13\textwidth}p{0.07\textwidth}p{0.11\textwidth}p{0.14\textwidth}p{0.18\textwidth}@{}}
  \caption{Material-parameter provenance.}\label{tab:provenance}\\
  \toprule
  Quantity & Symbol & Value/model & Unit & Temperature & Definition & Source \\
  \midrule
  \endfirsthead
  \toprule
  Quantity & Symbol & Value/model & Unit & Temperature & Definition & Source \\
  \midrule
  \endhead
  Band gap & $\Eg$ & \todo{ } & eV & \todo{ } & direct/indirect & \todo{cite} \\
  Electron mass & $\mathbf M_c$ & \todo{ } & $m_0$ & \todo{ } & curvature/DOS & \todo{cite} \\
  Hole mass & $\mathbf M_v$ & \todo{ } & $m_0$ & \todo{ } & HH/LH/DOS & \todo{cite} \\
  Valley count & $g_v$ & \todo{ } & 1 & \todo{ } & included in mass? & \todo{cite} \\
  Published $\Nc$ & $\Nc$ & \todo{ } & cm$^{-3}$ & \todo{ } & measured/derived & \todo{cite} \\
  Published $\Nv$ & $\Nv$ & \todo{ } & cm$^{-3}$ & \todo{ } & measured/derived & \todo{cite} \\
  Published $\nIntrinsic$ & $\nIntrinsic$ & \todo{ } & cm$^{-3}$ & \todo{ } & method & \todo{cite} \\
  \bottomrule
\end{longtable}
\normalsize

\subsection{Band-diagram interpretation}

\placeholderfigure{Published or properly redrawn/digitized band information.
Include attribution and reciprocal-space convention.}{fig:published-band}
\clearpage

\todo{Identify edge locations, direct/indirect character, multiplicity, nearby
bands, anisotropy, nonparabolicity, and limitations of the plotted path.}

\subsection{Material model and intrinsic results}

\todo{Define $\Eg(T)$, modeled components, temperature range, and any
material-specific use or extension of PhysKit.}

\placeholderfigure{Calculated $\Nc(T)$, $\Nv(T)$, $\mui(T)$, and/or $\nIntrinsic(T)$ for the selected material.}{fig:material-intrinsic}

\subsection{Published comparison and sensitivity}

\begin{table}[htbp]
  \centering
  \caption{Comparison with independent published data.}
  \label{tab:published-comparison}
  \begin{tabular}{@{}lllllll@{}}
    \toprule
    Quantity & $T$ & Calculated & Published & Unit & Difference & Source \\
    \midrule
    \todo{ } & \todo{ } & \todo{ } & \todo{ } & \todo{ } & \todo{ } & \todo{cite} \\
    \bottomrule
  \end{tabular}
\end{table}

\placeholderfigure{Sensitivity of $\nIntrinsic$ and $\mui$ to gap and electron/hole mass assumptions.}{fig:sensitivity}

\todo{Separate numerical error, parameter uncertainty, and physical model
limitations when interpreting disagreement.}

% =============================================================================
\section{Part IV: Equilibrium doping}
\label{sec:doping}
% =============================================================================

\subsection{Donor and acceptor models}

\begin{table}[htbp]
  \centering
  \scriptsize
  \caption{Dopant or defect-level provenance.}
  \label{tab:dopants}
  \begin{tabularx}{\textwidth}{@{}lXXXXXX@{}}
    \toprule
    Kind & Species/site & Level & Charge states & Degeneracy convention & Range & Source/limitations \\
    \midrule
    Donor & \todo{ } & \todo{ } & \todo{ } & \todo{ } & \todo{ } & \todo{cite} \\
    Acceptor & \todo{ } & \todo{ } & \todo{ } & \todo{ } & \todo{ } & \todo{cite} \\
    \bottomrule
  \end{tabularx}
\end{table}

With the explicitly declared degeneracy convention, one possible model is
\begin{align}
  N_D^+ &= \frac{N_D}{1+g_D\exp[(\mu-E_D)/(\kB T)]}, \\
  N_A^- &= \frac{N_A}{1+g_A\exp[(E_A-\mu)/(\kB T)]}, \\
  0 &= p+N_D^+-n-N_A^-.
\end{align}

\todo{Explain the convention, low/high-temperature behavior, solver bracket,
and neutrality residual.}

\subsection{Representative intrinsic, n-type, p-type, and compensated cases}

\begin{table}[htbp]
  \centering
  \caption{Modeled equilibrium cases.}
  \label{tab:doping-cases}
  \begin{tabular}{@{}lllll@{}}
    \toprule
    Case & $N_D$ & $N_A$ & Temperature range & Physical purpose \\
    \midrule
    Intrinsic & 0 & 0 & \todo{ } & baseline \\
    n-type & \todo{ } & \todo{ } & \todo{ } & \todo{ } \\
    compensated n-type & \todo{ } & \todo{ } & \todo{ } & \todo{ } \\
    p-type & \todo{ } & \todo{ } & \todo{ } & \todo{ } \\
    compensated p-type & \todo{ } & \todo{ } & \todo{ } & \todo{ } \\
    \bottomrule
  \end{tabular}
\end{table}

\placeholderfigure{$\mu(T)$ relative to band and dopant levels, with no implied spatial band bending.}{fig:doped-mu}
\placeholderfigure{Log-scale $n(T)$ and $p(T)$ for representative cases.}{fig:doped-carriers}
\placeholderfigure{Ionized fractions and neutrality residuals.}{fig:ionization}
\placeholderfigure{Carrier populations versus nominal dopant concentration.}{fig:concentration-sweep}

\subsection{Approximation comparison and regimes}

\todo{Identify freeze-out, extrinsic, intrinsic, and degenerate/non-Boltzmann
regimes where physically accessible. Compare incomplete and full ionization,
and state where the simpler result differs by less than five percent.}

% =============================================================================
\section{Synthesis and bounded validity}
\label{sec:synthesis}
% =============================================================================

\subsection{What transfers from Bloch2D to the real material?}

\todo{Separate transferable numerical ideas---curvature fitting, state counting,
occupation, convergence---from dimensionality and material-specific physics
that do not transfer directly.}

\subsection{Adequacy of the final model}

\todo{Answer quantitatively: over what temperature and concentration range is
the effective-mass plus equilibrium dopant model adequate, and which omitted
physical effect is most likely to invalidate it first?}

\subsection{Main conclusion}

\todo{Give a concise answer supported by your most decisive table and figure.}

% =============================================================================
\section{Reproducibility and validation record}
\label{sec:reproducibility}
% =============================================================================

\begin{table}[htbp]
  \centering
  \caption{Computational environment and execution record.}
  \label{tab:environment}
  \begin{tabularx}{\textwidth}{@{}p{0.28\textwidth}X@{}}
    \toprule
    Item & Value \\
    \midrule
    Repository/archive & \url{\repositoryurl} \\
    Python version & \todo{ } \\
    PhysKit version/commit & \physkitversion \\
    Main dependencies & \todo{ } \\
    Run order & \todo{ } \\
    Random seeds, if any & \todo{ } \\
    Clean-run result & \todo{ } \\
    Test/validation summary & \todo{ } \\
    \bottomrule
  \end{tabularx}
\end{table}

\todo{State how the submitted figures and tables can be regenerated. Disclose
materially reused code, data, and computational assistance according to course
policy.}

\clearpage
\bibliographystyle{unsrt}
\bibliography{references}

\appendix
\section{Supplementary convergence and validation evidence}

\todo{Place useful evidence here when it would interrupt the main argument.
Do not use the appendix to hide checks needed to support a main conclusion.}

\section{Additional parameter and provenance records}

\todo{Include extended tables, digitization details, or source-definition
comparisons if needed.}

\end{document}
